\emph{(1) Spin weights.} For a bivector $F$, $F^{2} = F\cdot F + F\wedge F$ lies in $\mathrm{span}\{1, I_{4}\}$, and with $F = \sum_{\mu<\nu}F^{\mu\nu}\gamma_{\mu}\gamma_{\nu}$, $\langle F^{2}\rangle_{0} = \sum_{\mu<\nu}(F^{\mu\nu})^{2}(\gamma_{\mu}\gamma_{\nu})^{2} = -\sum_{\mu<\nu}F^{\mu\nu}F_{\mu\nu} = -X/2$ (checked on three sample bivectors, one with $F\wedge F = 0$). The Lorentz generator of $F$ is $a \mapsto a\cdot F = \tfrac12(aF - Fa)$ on vectors (the derivative of $a \mapsto RaR^{\sim}$ at $R = 1$ along $\exp(-F/2)$) and $\psi \mapsto \tfrac12F\psi$ on spinors. The even algebra $\Cl^{+}(1,3)$ under left multiplication contains the minimal left ideal $L = \Cl^{+}(1,3)P$, $P = \tfrac12(1 + \gamma_{3}\gamma_{0})$, of real dimension four (checked); $I_{4}$ is central in the even algebra, so $L$ is $I_{4}$-stable (checked) and is a two-dimensional complex space with the unit $I_{4}$ of Definition \claimref{obj.field-theory.single-algebra-field}, the unit of Proposition \claimref{field-theory.second-quantisation}. Left multiplication by a bivector commutes with left multiplication by $I_{4}$, so $M(F)$ is complex-linear on $L$, and its complex trace is $\mathrm{tr}_{\C}T = \tfrac12\mathrm{Tr}_{\R}T - \tfrac12I_{4}\mathrm{Tr}_{\R}(I_{4}T)$. Exact traces give: on vectors $\mathrm{tr}\,M(F)^{2} = -X$, so $s = -4$; on $L$, $\mathrm{Re}\,\mathrm{tr}_{\C}M(F)^{2} = -X/4$, so $s = -1$; on the whole even algebra (eight real dimensions, two copies of $L$ under left multiplication) $s = -2$; on a scalar $s = 0$. On $L$ the spin coupling $E = 2M(F)$ is left multiplication by $F$, and $\mathrm{tr}_{\C}E^{2} = 2F^{2}$ in both its scalar and its $I_{4}$ part (checked); the $I_{4}$ part is nonzero exactly when $F\wedge F \neq 0$. Right multiplication by $F$ does not preserve $L$ (control). The whole even algebra is two copies of one chirality under left multiplication, which is the statement of Proposition \claimref{field-theory.second-quantisation}(4) that with $i = I_{4}$ the other chirality enters conjugated; its weight is the real part, the same for both chiralities, and the photon of Proposition \claimref{field-theory.fock-space-maxwell-modes} is the vector case.

\emph{(2) Weights per field.} Imported (Fact \claimref{synthesis.imported-heat-kernel-coefficients-and-background-field-gauge}): for a Laplace-type operator $-(D^{2} + E)$ on a flat background the logarithmic part of $\ln\det$ is $-\tfrac{1}{16\pi^{2}}\ln(L/\mu^{2})\int\mathrm{tr}\,a_{2}$ with $a_{2} = \tfrac{1}{12}\Omega_{\mu\nu}\Omega^{\mu\nu} + \tfrac12E^{2}$, $\Omega$ the curvature of $D$ (this is $a_{2}$ in DeWitt's numbering and $a_{4}$ in Vassilevich's, Phys. Rep. 388 (2003) 279, whose expansion runs in $t^{(k-n)/2}$, in his general form (4.15); its curvature terms vanish on the flat background); the background Feynman gauge is a covariant gauge, so this step is under the indefinite-form premise for the gauge field, Premise \claimref{synthesis.krein-form-premise-indefinite-inner-product-for-gauge-fields}, and the ghost subtraction is that of the gauge-fixing result (\claimref{synthesis.self-adjoint-brst-charge-needs-a-krein-form-and-the-tower-supplies-the-ghosts}, clause (iv)); in the background-field gauge (Abbott, Nucl. Phys. B185 (1981) 189) the one-loop operators are $-D^{2}$ for scalars and ghosts, the square of the Dirac operator (whose spin coupling is $E = $ left multiplication by the field strength, from $\gamma^{\mu}\gamma^{\nu}D_{\mu}D_{\nu} = D^{2} + \tfrac12\gamma^{\mu}\wedge\gamma^{\nu}[D_{\mu}, D_{\nu}]$) for spinors, and $-D^{2}\delta + 2\Omega$ acting on the vector index for the gauge field in Feynman gauge, with Faddeev-Popov ghosts a complex adjoint scalar of opposite statistics; the effective action is $\sigma w\ln\det$ with $w = 1$ for complex bosons and ghosts, $\tfrac12$ for real bosons and for the squared Dirac operator, and $\sigma = -1$ for fermions and ghosts. In each case $E = 2M(F)$ in the field's spin representation, tensored with the gauge generator; for gauge fields $F = F^{a}t^{a}$ and the internal trace gives the index, $\mathrm{tr}_{r}(t^{a}t^{b}) = T(r)\delta^{ab}$. So per unit index $\mathrm{tr}\,a_{2} \propto \tfrac{n}{12}X + \tfrac12\cdot 4\,\mathrm{tr}M(F)^{2} = \tfrac{X}{12}(n + 6s)$, the same overall factor for every field, and $b = \tfrac13\sigma w(n + 6s)$ in the normalisation of Fact \claimref{standard-model.imported-one-loop-running-trace-ratio} ($a_{i}(\mu) = a_{i}(M_{Z}) - \tfrac{b_{i}}{2\pi}\ln(\mu/M_{Z})$); the overall factor $\tfrac13$ is the standard one, and the Dirac value $\tfrac43$ it gives is the vacuum-polarisation coefficient $e^{2}/(12\pi^{2}) = \tfrac43\cdot e^{2}/(16\pi^{2})$ computed by Feynman parameters in the one-loop electrodynamics result (\claimref{synthesis.one-loop-qed-in-the-spacetime-algebra}). The spin weights entering here are those computed in (1), not inserted by hand (the check takes them from the traces). Then (exact): complex scalar $(\sigma, w, n, s) = (1, 1, 1, 0)$ gives $\tfrac13$; Weyl $(-1, \tfrac12, 2, -1)$ gives $\tfrac23$; Dirac $(-1, \tfrac12, 4, -2)$ gives $\tfrac43$, twice the Weyl value; the gauge vector $(1, \tfrac12, 4, -4)$ gives $-\tfrac{10}{3}$ and its ghosts $(-1, 1, 1, 0)$ give $-\tfrac13$, together $-\tfrac{11}{3}$. The orbital part (the $n$ term) of vector and ghosts is $\tfrac13(\tfrac12\cdot4 - 1) = \tfrac13$ and the spin part $\tfrac13\cdot\tfrac12\cdot6\cdot(-4) = -4$, the split of Nielsen and Hughes, $-(-1)^{2j}((2j)^{2} - \tfrac13)$ per helicity pair of spin $j$, which the check also evaluates.

\emph{(3) Indices.} Theorem \claimref{standard-model.gauge-algebra-minimal-massive-composite} gives $\su(2) \oplus \su(3) \oplus \uu(1)$ as the automorphisms of $\C^{2}\otimes\C^{3}$, acting on $\C^{2}$, $\C^{3}$ and on their exterior powers by derivations. With the hermitian generators $t$ (off-diagonal pairs and the diagonal ones, over $\mathbb{Q}(i)$), $T(r) = \tfrac12\,\mathrm{tr}(r(t)^{2})/\mathrm{tr}(t^{2})$ is the same for every generator and equals $\tfrac12$ on $\C^{2}$, $\C^{3}$, $\Lambda^{2}\C^{3}$ and $0$ on $\Lambda^{0}$, $\Lambda^{3}$; the adjoint Casimir, $\tfrac12\mathrm{tr}(\mathrm{ad}_{t}^{2})/\mathrm{tr}(t^{2})$ computed through the Gram matrix of the trace form, is $2$ and $3$ (all checked exactly). For $\su(2)$ the same ratio follows inside geometric algebra: in $\Cl(3,0)$ the pseudoscalar $e_{1}e_{2}e_{3}$ is central and squares to $-1$, so the even algebra is a two-dimensional complex space (a doublet) under left multiplication by bivectors, and for every bivector $B$ (checked on the three basis bivectors and one combination) $\mathrm{tr}(\mathrm{ad}_{B}^{2})$ on the bivectors is four times the complex trace of $L_{B}^{2}$ on the even algebra, which is $C_{2}(\mathrm{adj})/T(\C^{2}) = 2/\tfrac12$.

\emph{(4) Coefficients on the derived content.} By Theorem \claimref{standard-model.hypercharge-global-form-pair-counting-consistency}, under (S), the component $(j, k)$ carries $Y = (3j - 2k)/6$, the doublets of one generation are the even components $(1,1)$ and $(1,3)$, and the singlets are $(0, k)$ with the odd $k$ conjugated; as left-handed Weyl multiplets $(\dim_{2}, \Lambda^{k}, Y)$ this is $(2, \Lambda^{1}, \tfrac16)$, $(2, \Lambda^{3}, -\tfrac12)$, $(1, \Lambda^{0}, 0)$, $(1, \Lambda^{2}, \tfrac13)$, $(1, \Lambda^{2}, -\tfrac23)$, $(1, \Lambda^{0}, 1)$, recomputed from the rule (checked, with the cubic and gravitational hypercharge sums zero as a control). Summing $\tfrac23T$ over Weyl multiplets and $\tfrac13T$ over the complex components of $n_{h}$ doublets of hypercharge $\tfrac12$, and $-\tfrac{11}{3}C_{2}$ for each non-abelian factor, with $T_{1} = \tfrac35\sum Y^{2}$ (the normalisation $a_{1} = \tfrac35\alpha_{Y}^{-1}$): per generation $\sum T_{3} = 2$, $\sum T_{2} = 2$, $\sum\dim\,Y^{2} = \tfrac{10}{3}$, so $b_{3} = -11 + \tfrac23\cdot2n_{g}$, $b_{2} = -\tfrac{22}{3} + \tfrac23\cdot2n_{g} + \tfrac13\cdot\tfrac12n_{h}$, $b_{1} = \tfrac35(\tfrac23\cdot\tfrac{10}{3}n_{g} + \tfrac13\cdot\tfrac12n_{h})$, the stated values (checked for $0 \leq n_{g} \leq 9$, $0 \leq n_{h} \leq 4$, and as linear functions these determine them). Both counts are inputs: three generations is conjectured in the store and the doublet count is the rung field's, recorded as two. At $(3, 1)$ the coefficients are $(\tfrac{41}{10}, -\tfrac{19}{6}, -7)$, the values of Fact \claimref{standard-model.imported-one-loop-running-trace-ratio}; at $(3, 2)$ they are $(\tfrac{21}{5}, -3, -7)$; the supersymmetric listing is not reproduced (control). Finally $b_{3} < 0$ iff $4n_{g} < 33$, i.e. $n_{g} \leq 8$, and $b_{2} < 0$ iff $8n_{g} + n_{h} < 44$; at $n_{g} = 3$ the second condition is $n_{h} \leq 19$ (checked for $0 \leq n_{h} < 30$), so both non-abelian factors are asymptotically free whatever the doublet count up to nineteen; only $b_{1}$ and $b_{2}$ depend on it.

\emph{Scope.} The spin weights, indices and content are exact computations; the heat-kernel coefficient, the background-field gauge with its ghosts and the determinant weights are imported (Fact \claimref{synthesis.imported-heat-kernel-coefficients-and-background-field-gauge}), and the covariant gauge rests on the indefinite-form premise for the gauge field (Premise \claimref{synthesis.krein-form-premise-indefinite-inner-product-for-gauge-fields}). The result reproduces the stored import \claimref{standard-model.imported-coefficient-sum-indices} with its three coefficients computed rather than imported; it does not construct the non-abelian fields as quantum states.
